\documentclass[10pt,twocolumn]{article}

\usepackage[margin=0.75in,top=0.8in,bottom=0.8in]{geometry}
\usepackage{graphicx}
\usepackage{booktabs}
\usepackage{hyperref}
\usepackage{amsmath}
\usepackage{microtype}
\usepackage{caption}
\usepackage{float}
\usepackage{parskip}
\usepackage{titlesec}
\usepackage{enumitem}
\setlist{noitemsep, topsep=2pt}

\captionsetup{font=small, labelfont=bf}
\titlespacing{\section}{0pt}{6pt}{3pt}
\titlespacing{\subsection}{0pt}{4pt}{2pt}

\hypersetup{colorlinks=true, linkcolor=blue, citecolor=blue, urlcolor=blue}

\title{\textbf{Control-Plane Traffic in VxLAN Fabrics:\\BGP EVPN vs.\ OSPF Flood-and-Learn}\\[4pt]
\normalsize Course Research Project --- IP Routing and Protocols, April 2026}
\author{Ayush Mishra}
\date{}

\begin{document}
\maketitle
\thispagestyle{empty}

%% ============================================================
\begin{abstract}
Virtual Extensible LAN (VxLAN) standardises the data-plane encapsulation of
Ethernet frames over UDP, but leaves the choice of control plane---how VTEPs
learn remote MAC/IP bindings---entirely to the operator.  Two dominant
strategies exist: OSPF-underlaid \emph{flood-and-learn}, in which unknown
traffic is replicated to every VTEP in the fabric, and \emph{BGP EVPN}, which
distributes host reachability via MP-BGP before any data frame arrives.  Using
ContainerLab with FRRouting containers, we emulate spine-leaf fabrics at two
scales (4 and 8 VTEPs) and capture live packet data across ten runs per
scenario.  Our measurements show that per-host-addition control-plane traffic
in flood-and-learn is roughly \textbf{3.4--4.1$\times$} higher than in BGP
EVPN and grows closer to linearly with fabric size, while BGP EVPN grows
sub-linearly---providing concrete empirical grounding for the industry's
near-universal adoption of BGP EVPN in modern data centres.
\end{abstract}

%% ============================================================
\section{Introduction and Motivation}

Modern data centres require seamless Layer 2 connectivity across physically
distributed, Layer 3 routed infrastructure.  Tenant workloads must communicate
as if on the same LAN segment even when hosted on servers in different racks or
buildings, and east-west traffic now dominates the traffic matrix of
hyperscale providers such as AWS, Azure, and Google Cloud.

VxLAN (RFC 7348) addresses this by encapsulating Ethernet frames inside UDP
packets on port 4789, prepending a 24-bit VxLAN Network Identifier (VNI) that
supports over 16 million logical segments.  The boundary between the logical
overlay and the physical underlay is the \textbf{VTEP} (VxLAN Tunnel
Endpoint), typically a Top-of-Rack switch.  RFC 7348 deliberately leaves the
\emph{control plane}---how a VTEP learns which remote VTEP hosts a given
destination MAC---unspecified, creating the central design question this paper
investigates.

Without a control plane, VTEPs resort to \textbf{BUM flooding} (Broadcast,
Unknown unicast, Multicast): when a host sends to an unknown destination, the
originating VTEP replicates the frame and forwards a copy to \emph{every}
other VTEP in the fabric.  In a fabric with $n$ VTEPs, a single ARP request
generates $n-1$ replicated copies.  At 200 VTEPs with 100 hosts joining
simultaneously---a typical rack power-on scenario---this produces nearly
20,000 replicated packets within seconds, stressing uplink bandwidth and VTEP
CPU alike.

The practitioner's question is therefore: \emph{at what scale does the
overhead of flood-and-learn justify the added operational complexity of BGP
EVPN?}  Industry design guides advocate BGP EVPN qualitatively; this paper
provides side-by-side empirical measurements at fabric sizes representative of
small-to-medium deployments.

%% ============================================================
\section{VxLAN Background}

\subsection{What VxLAN Solves}

Traditional Ethernet is constrained by the 12-bit VLAN tag, limiting
deployments to 4,094 segments---far too few for multi-tenant clouds.  VxLAN's
24-bit VNI extends this to over 16 million segments.  More importantly, VxLAN
tunnels Layer 2 frames over routed IP infrastructure, allowing hypervisors on
physically distant servers to appear on the same broadcast domain.  A VTEP
receiving a native frame from a local host (1) looks up the destination MAC to
find the remote VTEP's IP, (2) wraps the frame in VxLAN/UDP/IP, and (3) sends
it across the underlay.  The remote VTEP strips the header and delivers the
original frame---transparent to the end hosts.

\subsection{The Control-Plane Gap}

RFC 7348 defines the encapsulation format but says nothing about how VTEPs
populate their MAC-to-VTEP forwarding tables.  Two strategies dominate:

\textbf{Flood-and-Learn} configures the underlay with a routing protocol (OSPF
in our experiments) to provide IP reachability between VTEPs.  The overlay
itself is stateless.  When a VTEP encounters an unknown destination it floods
the frame using \emph{ingress replication} (head-end replication): one
encapsulated copy per remote VTEP is sent individually.  VTEPs learn remote
MACs reactively by inspecting the source addresses of incoming encapsulated
frames---the same mechanism a conventional switch uses.

\textbf{BGP EVPN} (RFC 7432, RFC 8365) introduces a distributed control plane
on top of the same OSPF underlay.  Each VTEP acts as a BGP speaker and
advertises locally learned host bindings as \emph{EVPN Route Type 2
(MAC/IP Advertisement)} messages.  Spine switches act as BGP Route Reflectors,
redistributing each advertisement to all leaf VTEPs so that every VTEP knows
every host \emph{before} any data traffic arrives.  BGP EVPN also enables
\textbf{ARP suppression}: a VTEP that already holds a remote host's IP-to-MAC
binding can answer ARP requests locally, eliminating cross-fabric ARP floods
entirely.

%% ============================================================
\section{Implementation in Practice}

In a spine-leaf data centre fabric the BGP EVPN deployment follows a standard
pattern.  Spine switches run OSPF on their downlinks (the underlay) and also
act as iBGP Route Reflectors for the EVPN overlay.  Leaf switches (VTEPs) peer
with both spines via iBGP with the \texttt{l2vpn evpn} address family
activated.  Each leaf's VxLAN interface is created with the \texttt{nolearning}
flag set---disabling the flood-and-learn data plane---so that all forwarding
table population comes exclusively from BGP.

When a new host connects to Leaf 1, the sequence is: (1) the host sends a GARP
or ARP; (2) Leaf 1's kernel VXLAN driver detects the new MAC and FRRouting's
\texttt{bgpd} generates an EVPN Type 2 UPDATE; (3) the UPDATE propagates to
both spine Route Reflectors; (4) the spines reflect the UPDATE to all other
leaf peers; (5) each remote leaf installs the MAC/IP binding in its forwarding
table.  Steps 1--5 complete in seconds, \emph{before} any remote host needs to
reach the new host.  Subsequent traffic is sent directly to the correct VTEP
without any flooding.

In our experiments we used \textbf{ContainerLab} with \textbf{FRRouting (FRR)}
Docker containers as the emulation platform.  FRR provides production-grade
\texttt{ospfd} and \texttt{bgpd} daemons; VxLAN functionality is provided by
the Linux kernel's native VXLAN driver (the same code path used by Cumulus
Linux and SONiC).  Cisco NX-OSv/CSR images were excluded on resource grounds
(6--8 GB RAM per node vs.\ $\sim$150 MB for an FRR container).

%% ============================================================
\section{Experiments and Results}

\subsection{Setup and Metrics}

We emulated two spine-leaf topologies: a \textbf{small} fabric (2 spines + 4
leaf VTEPs) and a \textbf{medium} fabric (2 spines + 8 leaf VTEPs).  Each
scenario---flood-and-learn and BGP EVPN---was run ten times per fabric size, for
40 experimental runs total.  In each run a new host was added to Leaf 1 by
assigning an IP address and triggering ARP and ICMP traffic.  Packet captures
were collected with \texttt{tcpdump} on all leaf uplink interfaces and parsed
with Scapy to extract:

\begin{itemize}
  \item \textbf{Control-plane packet count}: VxLAN-encapsulated packets
        (flood-learn) or BGP UPDATE messages (BGP EVPN) per host-add event.
  \item \textbf{Control-plane byte volume}: aggregate bytes of the above.
  \item \textbf{Convergence time}: elapsed time from the first control-plane
        packet to the last in the capture window.
\end{itemize}

\subsection{Packet and Byte Overhead}

Figure~\ref{fig:bar_packets} shows the mean control-plane packet count per
host-add event across all four scenario-topology combinations; Figure~\ref{fig:bar_bytes}
shows the equivalent byte volumes.

\begin{figure}[H]
  \centering
  \includegraphics[width=\linewidth]{captures/plots/bar_packets.png}
  \caption{Mean control-plane packets per host-add event ($\pm$\,std\,dev,
  10 runs).  Blue bars: 4-leaf fabric; orange bars: 8-leaf fabric.
  Darker shades = BGP EVPN; lighter shades = flood-and-learn.}
  \label{fig:bar_packets}
\end{figure}

\begin{figure}[H]
  \centering
  \includegraphics[width=\linewidth]{captures/plots/bar_bytes.png}
  \caption{Mean control-plane byte volume per host-add event (KB,
  $\pm$\,std\,dev).  Same colour coding as Figure~\ref{fig:bar_packets}.}
  \label{fig:bar_bytes}
\end{figure}

In the 4-leaf fabric, flood-and-learn generated a mean of \textbf{86 VxLAN
packets} (10.3\,KB) per host-add event compared to \textbf{23 BGP UPDATE
packets} (2.5\,KB) for BGP EVPN---a \textbf{4.1$\times$ byte reduction}.  In
the 8-leaf fabric the gap widens: flood-and-learn produced \textbf{204 VxLAN
packets} (24.3\,KB) vs.\ \textbf{60 BGP UPDATE packets} (6.7\,KB) for BGP
EVPN---a \textbf{3.6$\times$ byte reduction} with substantially more absolute
overhead.  The BGP EVPN mean for the medium topology is skewed by the first
two runs in which BGP sessions were still being established; the steady-state
packet count (runs 4--10) was consistently \textbf{18 packets}---approximately
$1.8\times$ the 4-leaf steady-state of 10 packets, versus flood-and-learn
scaling by $2.4\times$ over the same doubling of VTEPs.

Table~\ref{tab:summary} summarises the mean values.

\begin{table}[H]
\centering
\caption{Summary of control-plane traffic per host-add event (non-zero runs).}
\label{tab:summary}
\small
\begin{tabular}{lrr}
\toprule
Scenario & Pkts & KB \\
\midrule
4-leaf Flood-Learn  & 86  & 10.3 \\
4-leaf BGP EVPN     & 23  & 2.5  \\
8-leaf Flood-Learn  & 204 & 24.3 \\
8-leaf BGP EVPN     & 60  & 6.7  \\
\bottomrule
\end{tabular}
\end{table}

\subsection{Scalability}

Figure~\ref{fig:scaling} plots mean control-plane packet count against fabric
size (number of VTEPs), clearly showing the diverging growth rates.

\begin{figure}[H]
  \centering
  \includegraphics[width=\linewidth]{captures/plots/scaling.png}
  \caption{Scalability: mean control-plane packets vs.\ number of VTEPs.
  Flood-and-learn (orange) grows faster than BGP EVPN (indigo) as the
  fabric expands.  Error bars show $\pm$\,std\,dev.}
  \label{fig:scaling}
\end{figure}

The flood-and-learn slope ($\approx$30 extra packets per additional VTEP) is
consistent with the ingress replication model: each additional VTEP is added to
every other VTEP's flood list, producing a near-linear increase.  BGP EVPN
grows more slowly ($\approx$9 packets per additional VTEP in steady state)
because the extra UPDATE messages are small, structured, and Route-Reflector
mediated---they do not require point-to-multipoint frame replication.

\subsection{Convergence Time}

Figure~\ref{fig:convergence} shows per-run convergence times.  Across both
scenarios and fabric sizes, convergence settled in the \textbf{6--10\,s} range.
Contrary to the hypothesis that BGP EVPN would converge faster, both protocols
showed similar median convergence times ($\approx$7.7\,s flood-learn vs.\
$\approx$8.1\,s BGP EVPN for the 8-leaf fabric).

\begin{figure}[H]
  \centering
  \includegraphics[width=\linewidth]{captures/plots/convergence.png}
  \caption{Convergence time (ms) per experimental run.  Solid lines: BGP EVPN;
  dashed lines: flood-and-learn.  Circles: 4-leaf fabric; squares: 8-leaf.}
  \label{fig:convergence}
\end{figure}

The similarity is attributable to the emulation environment: ContainerLab
virtual links and Docker networking introduce artificial latency that masks the
sub-second BGP propagation advantage seen on real hardware.  Specifically, BGP
session TCP setup and FRR daemon scheduling delays dominate the convergence
measurement, whereas on physical ASIC-based switches BGP EVPN convergence is
typically 100--500\,ms vs.\ seconds for flood-and-learn.  On physical gear,
flood-and-learn also suffers from the additional round-trip time of waiting for
the flooded ARP response before the forwarding table entry is installed.

The convergence measurements also exposed a warm-up artifact in the first one
or two runs of each experiment: BGP sessions were being established during the
run, inflating both the packet count and the convergence time.  Excluding these
warm-up runs (runs 1--2), steady-state measurements are highly consistent
across all remaining runs, validating the reproducibility of the emulated
testbed.

%% ============================================================
\section{Discussion and Conclusion}

Our experiments confirm, with empirical packet-level evidence, the qualitative
claims found in vendor design guides.  Flood-and-learn generates
\textbf{3.4--4.1$\times$ more control-plane traffic per host addition} than BGP
EVPN, and this advantage grows with fabric size.  At the 8-leaf scale the
absolute difference is already significant (204 vs.\ 60 packets; 24.3 vs.\
6.7\,KB per event); at hyperscale---hundreds of VTEPs---flood-and-learn becomes
operationally untenable.

The steady-state packet count for BGP EVPN ($\approx$10 at 4 leaves,
$\approx$18 at 8 leaves) reflects the bounded, protocol-defined cost of
propagating one EVPN Type 2 UPDATE through two Route Reflectors to all remote
VTEPs.  This cost is largely independent of the size of individual Ethernet
frames and does not flood tenant traffic onto unrelated VTEPs.

\textbf{Limitations.}  The emulated environment constrains convergence time
measurements and limits fabric scale to 8 VTEPs; physical hardware experiments
would strengthen the scalability argument.  Results also represent a single VNI
with low background traffic; production fabrics carry many VNIs and continuous
ARP/NDP background load that would widen the flood-and-learn overhead further.

\textbf{Design takeaway.}  For fabrics of 4 or fewer VTEPs, the 4.1$\times$
byte overhead of flood-and-learn may be acceptable given its simpler
configuration.  Beyond 8 VTEPs the overhead accumulates rapidly; BGP EVPN's
added configuration complexity pays for itself through substantially reduced
BUM traffic, eliminated ARP floods (via ARP suppression), and predictable
scaling behaviour.

%% ============================================================
\section*{References}
\begin{enumerate}[leftmargin=*, label={[\arabic*]}]
\small
  \item Mahalingam et~al., \emph{Virtual eXtensible Local Area Network
        (VXLAN)}, RFC 7348, IETF, 2014.
  \item Sajassi et~al., \emph{BGP MPLS-Based Ethernet VPN}, RFC 7432,
        IETF, 2015.
  \item Sajassi et~al., \emph{A Network Virtualization Overlay Solution
        Using EVPN}, RFC 8365, IETF, 2018.
  \item Cisco Systems, \emph{NX-OS VXLAN Configuration Guide for Nexus
        9000}, 2023.
  \item Arista Networks, \emph{EVPN Deployment and Design Guide}, 2022.
\end{enumerate}

\end{document}
